% !TEX root = ./Basilisk-opNavPoint-20180427.tex

\section{Test Description and Success Criteria}
The mathematics in this module are straight forward and can be tested in a series of input and output evaluation tests.


\subsection{Check 1}
Here a check is performed where the planet heading vector measurement $\bm h$ has a non-zero length and is not aligned with $\hat{\bm h}_{c}$.

\subsection{Check 2}
The planet direction vector $\bm h$ is given a norm value that is less than {\tt minUnitMag}.  In this case the attitude tracking $\bm\sigma_{B/R}$ should be set to zero.  Further, the body rate errors are now evaluated relative to a fixed $\bm\omega_{R/N}$ vector.

\subsection{Check 3}
The planet direction vector $\bm h$ aligned with $\hat{\bm h}_{c}$.  In this case the attitude tracking $\bm\sigma_{B/R}$ should be set to zero.  Further, the body rate errors are simply the inertial body angular rates.

\subsection{Check 4}
The planet direction vector is not given. The spacecraft is given a search rate as a sole command.

\subsection{Non-Zero Camera Mounting}
The camera frame is rotated relative to the body frame by several non-zero $\bm\sigma_{C/B}$ values and the spacecraft is placed at several attitudes $\bm\sigma_{B/N}$.  The reference attitude implied by the output, $[RN] = [BR]^{T}[BN]$, must point $\hat{\bm h}_{c}$ at the target.  The check is repeated after the measurement is dropped, so that the stored inertial heading is used.  The test {\tt test\_nonZeroCameraMounting} requires the boresight-to-target angle to be below $10^{-10}$ rad.

\subsection{Antiparallel Heading}
The target is placed directly behind $\hat{\bm h}_{c}$ for zero and non-zero $\bm\sigma_{C/B}$, with $\hat{\bm h}_{c}$ set to both $\hat{\bm c}_{3}$ and $\hat{\bm c}_{1}$ to exercise both cross products of the 180 degree fallback axis.  The test {\tt test\_antiparallelHeading} requires the boresight-to-target angle of the reference attitude to be below $10^{-7}$ rad.



\section{Test Parameters}
The unit test verify that the module output guidance message vectors match expected values.
\begin{table}[htbp]
	\caption{Error tolerance for each test.}
	\label{tab:errortol}
	\centering \fontsize{10}{10}\selectfont
	\begin{tabular}{ c | c } % Column formatting,
		\hline\hline
		\textbf{Output Value Tested}  & \textbf{Tolerated Error}  \\
		\hline
		$\bm\sigma_{B/R}$        & \input{AutoTex/toleranceValue}	   \\
		$\bm\omega_{B/R}$        & \input{AutoTex/toleranceValue}\\
		$\bm\omega_{R/N}$        & \input{AutoTex/toleranceValue} \\
		$\dot{\bm\omega}_{R/N}$        & \input{AutoTex/toleranceValue}  \\
		\hline\hline
	\end{tabular}
\end{table}

The nominal module test input values are $\hat{\bm h}_{c} = (0,0,1)$, $\bm h = (1,0,0)$ and $\leftexp{B}{\bm\omega}_{B/N} = (0.01, 0.50, -0.20)$ rad/sec.  The nominal body-fixed search rate is set to $\leftexp{B}{\bm\omega}_{R/N} = (0.0, 0.0, 0.1)$ rad/sec.  This rate is only used if no planet direction vector is available.  The small angle is set to $\epsilon$ = 0.01 degrees.



\section{Test Results}

All of the tests passed:
\begin{table}[H]
	\caption{Test results}
	\label{tab:results}
	\centering \fontsize{10}{10}\selectfont
	\begin{tabular}{c | c  } % Column formatting,
		\hline\hline
		\textbf{Check} 						  		&\textbf{Pass/Fail} \\
		\hline
	   1	   			& \input{AutoTex/passFail1} \\
	   2	   			& \input{AutoTex/passFail2} \\
	   3	   			& \input{AutoTex/passFail3} \\
	   4	   			& \input{AutoTex/passFail4} \\
	   \hline\hline
	\end{tabular}
\end{table}
